\chapter{Trabajos Relacionados} \label{cap:trabajos_relacionados}

Este capítulo revisa trabajos que aplican modelos de lenguaje al análisis de logs y a la observabilidad. Se agrupan por enfoque y se cierra con una comparación y la posición del presente trabajo. Las cifras provienen de los artículos tal como los resume el notebook de fuentes (\texttt{notes/04\_trabajos\_relacionados.md}); no son comparables directamente entre sí ni con las de esta tesis, pues usan conjuntos de datos y protocolos distintos. % TODO: verificar cada cifra contra el PDF original antes de la versión final.

\section{Revisiones sistemáticas}\label{sec:tr_revisiones}

He et al.~\cite{he2021_survey} y Jiang et al.~\cite{jiang2021_aiops} sintetizan el análisis automático de logs (\textit{parsing}, detección de anomalías, diagnóstico de fallas) y su papel en AIOps. Uddin et al.~\cite{uddin2026_microservice_logs} revisan específicamente el análisis de logs de microservicios con inteligencia artificial, y Wang y Qi~\cite{wang2024_rca} el análisis de causa raíz en (micro)servicios. Wang et al.~\cite{wang2025_slm_survey} y Chen et al.~\cite{chen2025_collab_survey} cubren los SLM y los mecanismos de colaboración entre modelos grandes y pequeños, respectivamente. Estas revisiones fundamentan el marco del Capítulo~\ref{cap:marco_teorico}.

\section{LLM con \textit{prompting} para análisis de logs}\label{sec:tr_prompting}

LogPrompt~\cite{liu2024_logprompt} propone estrategias de \textit{prompting} (\textit{self-prompt}, \textit{chain-of-thought} e \textit{in-context learning}) para analizar logs en línea sin entrenamiento, evaluado sobre ChatGPT y sobre un modelo abierto local (Vicuna-13B). Reporta un $F_1$ promedio de 0.797 en \textit{parsing} y valores de 0.384 (BGL) y 0.450 (Spirit) en detección de anomalías, además de explicaciones valoradas por expertos. Sus límites son la baja precisión absoluta en detección y la dependencia de la API en la nube para el mejor rendimiento.

Karlsen et al.~\cite{karlsen2024_benchmark} comparan LLM sobre tareas de análisis de logs, seguridad e interpretación. % TODO: describir resultados concretos tras consultar notes/04_trabajos_relacionados.md o el artículo (no se extrajeron cifras en esta redacción).

\section{Ajuste fino eficiente en parámetros para logs}\label{sec:tr_peft}

Lim et al.~\cite{lim2025_logadreft} comparan ReFT y LoRA sobre RoBERTa, GPT-2 y Llama-3 (8B) con los conjuntos HDFS, BGL, Spirit y Thunderbird. Llama3-ReFT alcanza $F_1$ de 0.9690 (BGL), 0.9950 (HDFS), 0.9192 (Spirit) y 0.9149 (Thunderbird), y supera a LoRA en 9 de 12 comparaciones; en generalización cruzada ReFT conserva un $F_1$ cercano a 0.50 mientras LoRA cae a valores cercanos a 0. El entrenamiento de ReFT es más lento que el de LoRA a rangos altos.

LINEADAPTER~\cite{liu2025_lineadapter} inserta un adaptador convolucional 1D tras la capa \textit{feed-forward} de bloques congelados (BERT, Llama-3.2-1B, Bloom-3B) y reporta $F_1$ de 0.988 (BGL), 0.980 (HDFS) y 0.953 (Hadoop) en detección, y similitud semántica de 0.959 a 0.977 con la causa raíz real, con unas cuatro veces menos parámetros y FLOPs entrenables que LoRA. Su evaluación se limita a un solo nodo de cómputo.

Estos trabajos usan conjuntos públicos de benchmark (HDFS, BGL, Spirit, Thunderbird) con etiquetas ya disponibles; el presente trabajo, en cambio, parte de un corpus real sin etiquetas y dedica un esfuerzo central a construirlas.

\section{Colaboración entre SLM y LLM}\label{sec:tr_colaboracion}

AdaptiveLog~\cite{ma2026_adaptivelog} combina un SLM local con un LLM: el SLM resuelve las muestras de baja incertidumbre y deriva las difíciles (alrededor del 27\,\%) al LLM, con lo que reduce los costos de API en un 73\,\% respecto de usar siempre el LLM. Reporta $F_1$ de 98.91\,\% en Thunderbird y 94.99\,\% en BGL. Mantiene la dependencia de un LLM externo o de servidores locales de gran capacidad.

\section{Observabilidad, RCA y detección con LLM en microservicios}\label{sec:tr_observabilidad}

MAIA~\cite{pedroso2025_maia} combina un Llama-3 8B ajustado de forma no supervisada, similitud coseno entre bloques de logs y una red bayesiana para RCA, y se evalúa en Kubernetes (Sock-Shop) con inyección de fallos mediante Chaos Mesh; los $F_1$ por tipo de fallo van de 0.83 a 0.997, con menor desempeño en fallos de E/S de base de datos no relacional. TRIDENT~\cite{dutta2025_trident} integra fusión multimodal de telemetría y ranking de causa raíz con atención sobre grafos; reporta reducciones de 25\,\% en MTTR y de 30\,\% en falsos positivos diarios, pero requiere instrumentación completa de trazas.
% TODO: otros trabajos del notebook sin entrada .bib (Wang et al. 2025 Preprints.org; TraceAI; Zhou et al. 2025; Mhatre y Shaikh 2024; Nellutla 2026; Awodele et al. 2026; LEMON). Agregar al .bib tras verificar si se desean citar.

\section{Modelos pequeños y despliegue local}\label{sec:tr_local}

Saha et al.~\cite{saha2025_slm_minimal} estudian el análisis de logs con SLM en entornos de recursos mínimos. SmartLLM~\cite{pardeshi2026_smartllm} propone un sistema \textit{on-premise} de modelos compactos con gobernanza y RAG, con latencia de 2 a 5 segundos y costo de API nulo, pero es una plataforma genérica sin especialización en observabilidad. Murphy et al.~\cite{murphy2024_embeddings} estudian \textit{embeddings} automáticos de mensajes de log, base del enfoque semántico usado aquí. % TODO: ampliar la descripción de Saha et al. con resultados concretos (modelos y métricas).

\section{Comparación y posición del presente trabajo}\label{sec:tr_comparacion}

\begin{table}[H]
\centering
\caption{Comparación cualitativa con trabajos relacionados. Los símbolos son una síntesis del autor a partir de lo descrito en las secciones anteriores.}
\label{tab:tr_comparacion}
\small
\begin{tabular}{@{}p{3.2cm}p{2.2cm}p{2.4cm}p{3.4cm}p{3.2cm}@{}}
\toprule
\textbf{Trabajo} & \textbf{Solo modelos locales} & \textbf{Datos} & \textbf{Método} & \textbf{Alerta y recomendación} \\
\midrule
LogPrompt~\cite{liu2024_logprompt} & No (API en la nube) & Benchmarks públicos & \textit{Prompting} & Explicaciones \\
AdaptiveLog~\cite{ma2026_adaptivelog} & No (híbrido) & Benchmarks públicos & SLM + LLM por incertidumbre & No \\
LogADReft~\cite{lim2025_logadreft} & Sí (ajuste local) & Benchmarks públicos & ReFT / LoRA & No \\
LINEADAPTER~\cite{liu2025_lineadapter} & Sí (ajuste local) & Benchmarks públicos & Adaptador Conv1D & RCA \\
MAIA~\cite{pedroso2025_maia} & Sí (Llama-3 8B) & Sock-Shop (sintético) & Autorregresivo + Bayes & RCA \\
\textbf{Esta tesis} & Sí (Ollama / HF) & Corpus real DTIC & SLM + embeddings, \textit{prompt} y LoRA & Slack, recomendador, \textit{dashboard} \\
\bottomrule
\end{tabular}
\normalsize
\end{table}

% TODO: validar la tabla (columnas "Solo modelos locales" y "Datos") contra los artículos originales; la clasificación proviene del resumen del notebook.

Frente a la literatura, el presente trabajo se diferencia en cuatro puntos: (i) trabaja sobre un corpus operativo real de la universidad, sin etiquetas previas, y documenta cómo se construyó el \textit{ground truth}; (ii) compara varios SLM de distintas familias y tamaños con la misma tarea, sin depender de servicios en la nube; (iii) evalúa de forma semántica con \textit{embeddings} y permite comparar tamaños de \textit{batch}; (iv) integra alerta, recomendación e inyección de errores en un mismo \textit{pipeline} con interfaz gráfica.
